% Lecture 7 notes: \input by notes.tex, \subimport'ed by the semester book.
% Exam questions (экзамен): 16 (критерий Эйлера, метод Чиполлы), 17 (лемма Гензеля, сравнения
% по модулю p^k).
\begin{lecture}{7}{Квадратичные сравнения. Методы Чиполлы и Гензеля}

\section{Зачем извлекать квадратные корни по модулю}

Задача этой лекции "--- решить сравнение
\[
  x^2 \equiv a \Mod{m},
\]
т.\,е. извлечь квадратный корень из $a$ в кольце $\Zn{m}$. В криптографии она возникает
постоянно.
\begin{itemize}
  \item \emph{Эллиптические кривые} (лекция~8). Точка кривой $y^2 = x^3 + Ax + B$ над $\GF{p}$
    хранится и передаётся «в сжатом виде»: координата $x$ и один бит $y$. Чтобы восстановить
    $y$, нужно извлечь корень из $x^3 + Ax + B$ по модулю $p$. Так же устроено отображение
    сообщения или хеш-значения в точку кривой.
  \item \emph{Квадратичное решето} (лекция~6). Чтобы просеивать значения $Q(t) = t^2 - N$,
    нужно для каждого простого $p$ факторной базы найти корни сравнения
    $t^2 \equiv N \Mod{p}$, а для учёта степеней простых "--- корни по модулю $p^k$.
  \item \emph{Криптосистема Рабина} и протоколы, основанные на сложности задачи о квадратичных
    вычетах (шифрование Гольдвассер"--~Микали). Для модуля $N = pq$ с неизвестным разложением
    извлечение корней оказывается \emph{равносильным факторизации} $N$
    (теорема~\ref{l7:rabin}), а для простого модуля это лёгкая задача.
\end{itemize}

План лекции: сначала "--- простой модуль $p$ (критерий Эйлера, простые случаи, алгоритмы
Тоннелли"--~Шенкса и Чиполлы), затем "--- степени простого (лемма Гензеля), наконец "---
составной модуль (китайская теорема об остатках) и связь с факторизацией.

Всюду в разделах~\ref{l7:sec-qr}--\ref{l7:sec-cipolla} $p$ "--- нечётное простое число.

\section{Квадратичные вычеты и критерий Эйлера}\label{l7:sec-qr}

\begin{definition}
Число $a$, $p \nmid a$, называется \emph{квадратичным вычетом} по модулю $p$, если сравнение
$x^2 \equiv a \Mod{p}$ разрешимо, и \emph{квадратичным невычетом} в противном случае.
\end{definition}

Число $0$ мы не относим ни к вычетам, ни к невычетам: сравнение $x^2 \equiv 0 \Mod{p}$ имеет
единственное решение $x \equiv 0$.

\begin{proposition}\label{l7:count}
Если $x^2 \equiv a \Mod{p}$ разрешимо и $p \nmid a$, то решений ровно два: $\pm x_0$.
Среди чисел $1, 2, \dots, p - 1$ ровно $\frac{p-1}{2}$ квадратичных вычетов и
$\frac{p-1}{2}$ невычетов; вычеты "--- это числа $1^2, 2^2, \dots, \bigl(\frac{p-1}{2}\bigr)^2$
по модулю $p$.
\end{proposition}

\begin{proof}
Если $x_0^2 \equiv x^2 \equiv a$, то $p \mid (x - x_0)(x + x_0)$, и так как $p$ простое,
$x \equiv \pm x_0$. При этом $x_0 \not\equiv -x_0$, иначе $p \mid 2x_0$, что невозможно для
нечётного $p$ и $p \nmid x_0$. Значит, отображение $x \mapsto x^2$ группы $\Znx{p}$ в себя
«склеивает» элементы ровно парами $\{x, -x\}$, и его образ (множество вычетов) состоит из
$\frac{p-1}{2}$ элементов. Каждая пара $\{x, -x\}$ содержит ровно одно число из
$1, \dots, \frac{p-1}{2}$, откуда последнее утверждение.
\end{proof}

\begin{example}
По модулю $11$: $1^2 = 1$, $2^2 = 4$, $3^2 = 9$, $4^2 = 16 \equiv 5$, $5^2 = 25 \equiv 3$.
Вычеты "--- $\{1, 3, 4, 5, 9\}$, невычеты "--- $\{2, 6, 7, 8, 10\}$.
\end{example}

Для больших $p$ перебор квадратов бесполезен. Быстрый способ проверить разрешимость даёт

\begin{theorem}[критерий Эйлера]\label{l7:euler}
Пусть $p \nmid a$. Тогда
\[
  a^{(p-1)/2} \equiv \begin{cases}
    \phantom{-}1 \Mod{p}, & a \text{ "--- квадратичный вычет},\\
    -1 \Mod{p}, & a \text{ "--- квадратичный невычет}.
  \end{cases}
\]
\end{theorem}

\begin{proof}
Обозначим $u = a^{(p-1)/2}$. По малой теореме Ферма $u^2 = a^{p-1} \equiv 1$, поэтому
$p \mid (u - 1)(u + 1)$ и $u \equiv \pm 1$.

Если $a \equiv x^2$ "--- вычет, то $u \equiv x^{p-1} \equiv 1$. Значит, все $\frac{p-1}{2}$
вычетов (утверждение~\ref{l7:count}) являются корнями многочлена
$X^{(p-1)/2} - 1$ над полем $\GF{p}$. Многочлен степени $\frac{p-1}{2}$ над полем имеет не более
$\frac{p-1}{2}$ корней, поэтому других корней у него нет: для невычета $a$ имеем
$u \not\equiv 1$, т.\,е. $u \equiv -1$.
\end{proof}

Проверка по критерию Эйлера "--- это одно возведение в степень, т.\,е. $\bigO(\log p)$ умножений
по модулю $p$ (лекция~6 первого семестра).

\begin{remark}
Другое доказательство использует первообразный корень $g$ (лекция~7 первого семестра): если
$a = g^k$, то $a$ "--- вычет тогда и только тогда, когда $k$ чётно, а
$a^{(p-1)/2} = g^{k(p-1)/2} \equiv 1$ тогда и только тогда, когда $(p - 1) \mid k\frac{p-1}{2}$,
т.\,е. снова когда $k$ чётно.
\end{remark}

\subsection{Символ Лежандра}

\begin{definition}
\emph{Символ Лежандра} числа $a$ по простому модулю $p$:
\[
  \legendre{a}{p} = \begin{cases}
    0, & p \mid a,\\
    1, & a \text{ "--- квадратичный вычет},\\
    -1, & a \text{ "--- квадратичный невычет}.
  \end{cases}
\]
\end{definition}

Критерий Эйлера записывается так: $\legendre{a}{p} \equiv a^{(p-1)/2} \Mod{p}$, а число решений
сравнения $x^2 \equiv a \Mod{p}$ равно $1 + \legendre{a}{p}$.

\begin{proposition}\label{l7:legendre}
\begin{enumerate}
  \item Если $a \equiv b \Mod{p}$, то $\legendre{a}{p} = \legendre{b}{p}$.
  \item $\legendre{ab}{p} = \legendre{a}{p}\legendre{b}{p}$ (мультипликативность);
    в частности, $\legendre{a b^2}{p} = \legendre{a}{p}$ при $p \nmid b$.
  \item $\legendre{-1}{p} = (-1)^{(p-1)/2}$: $-1$ "--- вычет тогда и только тогда, когда
    $p \equiv 1 \Mod 4$.
  \item $\legendre{2}{p} = (-1)^{(p^2-1)/8}$: $2$ "--- вычет тогда и только тогда, когда
    $p \equiv \pm 1 \Mod 8$.
  \item (\emph{квадратичный закон взаимности}) Для различных нечётных простых $p, q$
    \[
      \legendre{p}{q}\legendre{q}{p} = (-1)^{\frac{p-1}{2}\cdot\frac{q-1}{2}}.
    \]
\end{enumerate}
\end{proposition}

\begin{proof}
Свойство~1 очевидно. Свойства~2 и~3 следуют из критерия Эйлера: обе части "--- числа $\pm 1$
(или $0$), сравнимые по модулю $p > 2$, а значит, равные:
$(ab)^{(p-1)/2} = a^{(p-1)/2} b^{(p-1)/2}$. Свойства~4 и~5 мы примем без доказательства
(их доказательства с помощью леммы Гаусса см. в~\cite[гл.~5]{ireland-rosen},
\cite[гл.~V]{vinogradov}).
\end{proof}

Свойства~1--5 позволяют вычислять символ Лежандра, не возводя в степень, "--- так же, как
алгоритм Евклида, с помощью символа Якоби (лекция~8 первого семестра, тест
Соловея"--~Штрассена). Например, $\legendre{10}{89} = \legendre{2}{89}\legendre{5}{89} =
1 \cdot \legendre{89}{5} = \legendre{4}{5} = 1$, так как $89 \equiv 1 \Mod 8$ и
$5 \equiv 1 \Mod 4$.

\checkyourself{Почему критерий Эйлера сразу даёт мультипликативность символа Лежандра? Является ли
произведение двух невычетов вычетом?}

\section{\texorpdfstring{Простые случаи: $p \equiv 3 \pmod 4$ и $p \equiv 5 \pmod 8$}{Простые случаи: p = 3 (mod 4) и p = 5 (mod 8)}}\label{l7:sec-easy}

Критерий Эйлера отвечает на вопрос, \emph{есть ли} корень. Найти его в половине случаев можно
одной формулой.

\begin{theorem}\label{l7:p34}
Пусть $p \equiv 3 \Mod 4$ и $a$ "--- квадратичный вычет. Тогда $x = a^{(p+1)/4} \bmod p$ "---
корень сравнения $x^2 \equiv a \Mod{p}$.
\end{theorem}

\begin{proof}
Показатель $\frac{p+1}{4}$ целый. По критерию Эйлера
$x^2 = a^{(p+1)/2} = a \cdot a^{(p-1)/2} \equiv a \cdot 1 = a$.
\end{proof}

Если $a$ "--- невычет, то та же выкладка даёт $x^2 \equiv -a$; поэтому можно сразу вычислить
$x = a^{(p+1)/4}$ и проверить, равно ли $x^2$ числу $a$: если нет "--- решений нет.

\begin{example}
$p = 23 \equiv 3 \Mod 4$, $a = 2$: $x = 2^{6} = 64 \equiv 18$, $18^2 = 324 \equiv 2 \Mod{23}$.
Корни: $x \equiv \pm 18$, т.\,е. $\{5, 18\}$.
\end{example}

\taskexample{sqrt-p34}

\begin{proposition}[формула Аткина]\label{l7:atkin}
Пусть $p \equiv 5 \Mod 8$ и $a$ "--- квадратичный вычет. Положим
\[
  b = (2a)^{(p-5)/8}, \qquad i = 2ab^2, \qquad x = ab(i - 1) \quad(\bmod p).
\]
Тогда $x^2 \equiv a \Mod{p}$.
\end{proposition}

\begin{proof}
Имеем $i = 2a \cdot (2a)^{(p-5)/4} = (2a)^{(p-1)/4}$. По критерию Эйлера и
свойству~4 утверждения~\ref{l7:legendre} $i^2 = (2a)^{(p-1)/2} \equiv
\legendre{2}{p}\legendre{a}{p} = -1 \cdot 1 = -1$. Тогда
$x^2 = a^2 b^2 (i - 1)^2 = a^2 b^2 (i^2 - 2i + 1) = a^2 b^2 \cdot (-2i) = -a \cdot i \cdot 2ab^2
= -a i^2 = a$.
\end{proof}

\begin{example}
$p = 29 \equiv 5 \Mod 8$, $a = 5$ (вычет: $\legendre{5}{29} = \legendre{29}{5} = \legendre{4}{5} = 1$).
$b = 10^{3} = 1000 \equiv 14$, $i = 10 \cdot 14^2 = 1960 \equiv 17$ (действительно,
$17^2 = 289 \equiv -1$), $x = 5 \cdot 14 \cdot 16 = 1120 \equiv 18$; проверка:
$18^2 = 324 \equiv 5 \Mod{29}$.
\end{example}

Остаётся случай $p \equiv 1 \Mod 8$, и он принципиально труднее: все известные быстрые
алгоритмы используют какой-нибудь \emph{квадратичный невычет}. Невычетов ровно половина, и
случайное число оказывается невычетом с вероятностью $\frac12$, но детерминированного
полиномиального алгоритма поиска невычета не известно (в предположении обобщённой гипотезы
Римана наименьший невычет меньше $2\ln^2 p$). Поэтому алгоритмы этого раздела
\emph{вероятностные} типа Лас-Вегас: ответ всегда верен, случайно лишь время работы.

\checkyourself{Почему формула $x = a^{(p+1)/4}$ не годится при $p \equiv 1 \Mod 4$? Что получится,
если применить её к невычету $a$ при $p \equiv 3 \Mod 4$?}

\section{Алгоритм Тоннелли"--~Шенкса (кратко)}\label{l7:sec-ts}

Запишем $p - 1 = 2^s q$, $q$ нечётно. Группа $\Znx{p}$ циклическая, её силовская
$2$-подгруппа $S$ (элементы порядка, делящего $2^s$) "--- циклическая порядка $2^s$, и её
порождает $c = z^q$ для любого невычета $z$ (порядок $c$ равен $2^s$, так как
$c^{2^{s-1}} = z^{(p-1)/2} = -1$).

\begin{algorithm}[Тоннелли"--~Шенкс]
\begin{algorithmic}[1]
\Require нечётное простое $p = 2^s q + 1$, квадратичный вычет $a$.
\Ensure $x$: $x^2 \equiv a \Mod p$.
\State найти невычет $z$ (случайный выбор и проверка критерием Эйлера)
\State $c \gets z^q$, $x \gets a^{(q+1)/2}$, $b \gets a^q$, $m \gets s$
\While{$b \ne 1$}
  \State найти наименьшее $i \ge 1$ с $b^{2^i} = 1$ \Comment{$i < m$}
  \State $d \gets c^{2^{m - i - 1}}$, $x \gets x d$, $c \gets d^2$, $b \gets b d^2$, $m \gets i$
\EndWhile
\State \Return $x$
\end{algorithmic}
\end{algorithm}

На каждом шаге сохраняются инварианты $x^2 = ab$, $\ord c = 2^m$, $\ord b$ делит $2^{m-1}$
(вначале $x^2 = a^{q+1} = a b$, а $b = a^q$ лежит в $S$ и является квадратом, поэтому его
порядок меньше $2^s$). Если $\ord b = 2^i$, то $d^2 = c^{2^{m-i}}$ имеет порядок ровно $2^i$, и
$b d^2$ имеет порядок \emph{строго меньше} $2^i$ (оба элемента порядка $2^i$ в циклической
группе $S$, и $(bd^2)^{2^{i-1}} = (-1)(-1) = 1$). Порядок $b$ убывает, и не более чем за $s$
итераций $b = 1$, т.\,е. $x^2 = a$. Сложность "--- $\bigO(\log p + s^2)$ умножений по модулю $p$
плюс поиск невычета. При $p \equiv 3 \Mod 4$ ($s = 1$) цикл не выполняется ни разу и алгоритм
совпадает с теоремой~\ref{l7:p34}. Неудобный случай "--- большое $s$ (например,
$p = 2^{s} q + 1$ с $q$ маленьким): тогда до $s^2 \sim \log^2 p$ умножений. Алгоритм Чиполлы
лишён этого недостатка.

\section{Метод Чиполлы}\label{l7:sec-cipolla}

Идея: если $\legendre{a}{p} = 1$, то многочлен $X^2 - a$ раскладывается на линейные
множители над $\GF{p}$; найдём его корень, перейдя в квадратичное расширение $\GF{p^2}$, где
вычисления «видят» сопряжение.

\subsection{\texorpdfstring{Шаг 1: подходящее $t$}{Шаг 1: подходящее t}}

\begin{lemma}\label{l7:cipolla-t}
Пусть $a$ "--- квадратичный вычет по модулю $p$. Ровно для $\frac{p-1}{2}$ значений
$t \in \GF{p}$ число $t^2 - a$ является квадратичным невычетом.
\end{lemma}

\begin{proof}
Посчитаем $t$, для которых $t^2 - a$ "--- \emph{ненулевой квадрат}, т.\,е.
$t^2 - a = s^2$ при некотором $s \ne 0$. Равенство $t^2 - s^2 = a$ перепишем как $uv = a$, где
$u = t - s$, $v = t + s$. Замена $(t, s) \leftrightarrow (u, v)$ "--- биекция $\GF{p}^2 \to
\GF{p}^2$ (обратная: $t = \frac{u+v}{2}$, $s = \frac{v-u}{2}$; здесь важно, что $p$ нечётно).
Уравнение $uv = a$ имеет $p - 1$ решений ($u \ne 0$ любое, $v = a/u$). Из них решений с
$s = 0$, т.\,е. $u = v$, $u^2 = a$, ровно два. Остаётся $p - 3$ пар $(t, s)$ с $s \ne 0$;
каждому подходящему $t$ соответствуют две пары $(t, \pm s)$, т.\,е. таких $t$ ровно
$\frac{p-3}{2}$.

Далее, $t^2 - a = 0$ ровно для двух значений $t$ (корней из $a$). Все остальные
\[
  p - \frac{p-3}{2} - 2 = \frac{p - 1}{2}
\]
значений $t$ дают невычеты.
\end{proof}

Таким образом, случайное $t$ подходит с вероятностью $\frac{p-1}{2p} \approx \frac12$, и
среднее число испытаний (каждое "--- проверка $t^2 - a$ критерием Эйлера) равно
$\frac{2p}{p-1} \approx 2$. Чаще перебирают $t = 1, 2, 3, \dots$ подряд; на практике это тоже
работает быстро, хотя строгой оценки для такого перебора нет.

\subsection{\texorpdfstring{Шаг 2: вычисления в $\GF{p^2}$}{Шаг 2: вычисления в F(p\textasciicircum 2)}}

Пусть $d = t^2 - a$ "--- невычет. Многочлен $X^2 - d$ не имеет корней в $\GF{p}$, значит,
неприводим (у приводимого многочлена степени~$2$ есть линейный множитель), и
\[
  \GF{p^2} = \GF{p}[X]/(X^2 - d) = \{u + v\omega : u, v \in \GF{p}\}, \qquad \omega^2 = d,
\]
"--- поле из $p^2$ элементов.
Умножение выполняется по правилу
\[
  (u_1 + v_1\omega)(u_2 + v_2\omega) = (u_1u_2 + d\,v_1v_2) + (u_1v_2 + u_2v_1)\,\omega,
\]
т.\,е. стоит пяти умножений в $\GF{p}$ (четырёх, если $d$ мало, и меньше "--- с приёмом
Карацубы).

\begin{lemma}[автоморфизм Фробениуса]\label{l7:frob}
Для всех $u, v \in \GF{p}$ в поле $\GF{p}(\omega)$ выполнено $(u + v\omega)^p = u - v\omega$.
\end{lemma}

\begin{proof}
В кольце характеристики $p$ верно $(\alpha + \beta)^p = \alpha^p + \beta^p$: биномиальные
коэффициенты $\binom{p}{k}$ при $0 < k < p$ делятся на $p$. Поэтому
$(u + v\omega)^p = u^p + v^p \omega^p$. По малой теореме Ферма $u^p = u$, $v^p = v$, а
$\omega^p = \omega \cdot (\omega^2)^{(p-1)/2} = \omega\, d^{(p-1)/2} = -\omega$ по критерию
Эйлера, так как $d$ "--- невычет.
\end{proof}

\begin{theorem}[Чиполла]\label{l7:cipolla}
Пусть $a$ "--- квадратичный вычет, $d = t^2 - a$ "--- невычет, $\omega^2 = d$. Тогда
\[
  x = (t + \omega)^{(p+1)/2} \in \GF{p}, \qquad x^2 \equiv a \Mod{p}.
\]
\end{theorem}

\begin{proof}
По лемме~\ref{l7:frob}
\[
  x^2 = (t + \omega)^{p+1} = (t + \omega)^p (t + \omega) = (t - \omega)(t + \omega) = t^2 - \omega^2
  = t^2 - d = a.
\]
Осталось показать, что $x$ лежит в $\GF{p}$ (т.\,е. коэффициент при $\omega$ равен нулю). Так как
$a$ "--- вычет, у многочлена $X^2 - a$ есть корни $\pm r \in \GF{p}$. Над полем $\GF{p^2}$ у
многочлена степени~$2$ не более двух корней, а $x$ "--- его корень, поэтому $x \in \{r, -r\}
\subset \GF{p}$.
\end{proof}

\begin{algorithm}[Чиполла]
\begin{algorithmic}[1]
\Require нечётное простое $p$, $a$ с $\legendre{a}{p} = 1$.
\Ensure $x$: $x^2 \equiv a \Mod p$.
\Repeat
  \State выбрать $t \in \GF{p}$ (случайно или $t = 1, 2, \dots$), $d \gets t^2 - a$
\Until{$d^{(p-1)/2} \equiv -1$} \Comment{в среднем $\approx 2$ попытки}
\State вычислить $x + y\omega \gets (t + \omega)^{(p+1)/2}$ бинарным алгоритмом в $\GF{p}(\omega)$,
  $\omega^2 = d$
\State \Return $x$ \Comment{$y = 0$}
\end{algorithmic}
\end{algorithm}

\emph{Сложность.} Каждая попытка "--- одно возведение в степень $\frac{p-1}{2}$ в $\GF{p}$, т.\,е.
$\bigO(\log p)$ умножений; в среднем попыток две. Возведение в степень $\frac{p+1}{2}$ в
$\GF{p^2}$ "--- $\bigO(\log p)$ умножений в $\GF{p^2}$, каждое "--- $\bigO(1)$ умножений в $\GF{p}$.
Итого $\bigO(\log p)$ умножений по модулю $p$, т.\,е. $\bigO(\log^3 p)$ битовых операций
школьным умножением, \emph{независимо} от степени двойки в $p - 1$.

\begin{example}\label{l7:ex-cipolla}
$p = 13$, $a = 10$; $10^6 = (10^2)^3 = 9^3 = 729 \equiv 1$, т.\,е. $a$ "--- вычет. Подбор $t$:
$t = 1$: $1 - 10 \equiv 4 = 2^2$ "--- вычет; $t = 2$: $4 - 10 \equiv 7$, $7^6 \equiv
(7^2)^3 = 10^3 \equiv 12 \equiv -1$ "--- невычет. Итак, $\omega^2 = 7$, вычисляем
$(2 + \omega)^7$, $7 = 111_2$:
\[
  \begin{array}{c|c|c}
    \text{бит} & r^2 & r \\ \hline
    1 & 1 & 2 + \omega \\
    1 & (2 + \omega)^2 = 11 + 4\omega & (11 + 4\omega)(2 + \omega) = 11 + 6\omega \\
    1 & (11 + 6\omega)^2 = 9 + 2\omega & (9 + 2\omega)(2 + \omega) = 6
  \end{array}
\]
(например, $(11 + 4\omega)(2 + \omega) = 22 + 4 \cdot 7 + (11 + 8)\omega = 50 + 19\omega \equiv
11 + 6\omega$). Получили $x = 6$: $6^2 = 36 \equiv 10 \Mod{13}$; корни $\pm 6$, т.\,е.
$\{6, 7\}$.
\end{example}

\taskexample{sqrt-cipolla}

\checkyourself{Где в доказательстве теоремы Чиполлы используется то, что $t^2 - a$ "--- именно
невычет? Что вычислит алгоритм, если по ошибке взять $t$, для которого $t^2 - a$ "--- вычет?}

\checkyourself{Что произойдёт в алгоритме Чиполлы, если $a$ "--- невычет? Почему перед запуском
нужно проверить критерий Эйлера?}

\section{Лемма Гензеля}\label{l7:sec-hensel}

Перейдём к модулю $p^k$. Идея та же, что у метода Ньютона для вещественных корней: корень по
модулю $p$ «уточняется» до корня по модулю $p^2, p^3, \dots$ с помощью производной. Ниже
$f \in \ZZ[X]$ "--- многочлен с целыми коэффициентами, $f'$ "--- его (формальная) производная,
$p$ "--- любое простое.

\begin{lemma}\label{l7:taylor}
Для любых целых $x, h$ выполнено $f(x + h) = f(x) + h f'(x) + h^2 g(x, h)$, где
$g \in \ZZ[X, H]$.
\end{lemma}

\begin{proof}
Достаточно проверить для $f = X^n$: $(x + h)^n = x^n + n x^{n-1} h + \sum_{j \ge 2}
\binom{n}{j} x^{n-j} h^j$, и все биномиальные коэффициенты целые. Общий случай получается
линейностью.
\end{proof}

\begin{theorem}[лемма Гензеля]\label{l7:hensel}
Пусть $k \ge 1$, $f(x_k) \equiv 0 \Mod{p^k}$ и $f'(x_k) \not\equiv 0 \Mod{p}$. Тогда существует
единственный по модулю $p^{k+1}$ корень $x_{k+1}$ сравнения $f(x) \equiv 0 \Mod{p^{k+1}}$ с
$x_{k+1} \equiv x_k \Mod{p^k}$. Он равен
\[
  x_{k+1} = x_k + t p^k, \qquad t \equiv -\frac{f(x_k)}{p^k} \cdot f'(x_k)^{-1} \Mod{p},
\]
или, что то же самое, $x_{k+1} \equiv x_k - f(x_k)\, f'(x_k)^{-1} \Mod{p^{k+1}}$.
\end{theorem}

\begin{proof}
Любое число, сравнимое с $x_k$ по модулю $p^k$, имеет вид $x_k + tp^k$, и его остаток по модулю
$p^{k+1}$ определяется остатком $t \bmod p$. По лемме~\ref{l7:taylor} с $h = tp^k$
\[
  f(x_k + tp^k) = f(x_k) + t p^k f'(x_k) + t^2 p^{2k} g \equiv f(x_k) + t p^k f'(x_k)
  \Mod{p^{k+1}},
\]
так как $2k \ge k + 1$. По условию $f(x_k) = p^k s$ с целым $s$. Тогда
\[
  f(x_k + tp^k) \equiv p^k \bigl(s + t f'(x_k)\bigr) \Mod{p^{k+1}},
\]
и это $\equiv 0$ тогда и только тогда, когда $s + t f'(x_k) \equiv 0 \Mod{p}$. Так как
$f'(x_k)$ обратим по модулю $p$, последнее сравнение имеет единственное решение
$t \equiv -s f'(x_k)^{-1} \Mod p$.

Вторая формула: $x_k - f(x_k) f'(x_k)^{-1} = x_k - p^k s\, f'(x_k)^{-1}$, а по модулю
$p^{k+1}$ слагаемое $p^k \cdot s f'(x_k)^{-1}$ зависит лишь от $s f'(x_k)^{-1} \bmod p$.
\end{proof}

Заметим, что $f'(x_k) \equiv f'(x_1) \Mod p$ (так как $x_k \equiv x_1 \Mod p$), поэтому обратный
элемент $f'(x_1)^{-1} \bmod p$ достаточно найти \emph{один раз}. Применяя теорему
последовательно при $k = 1, 2, \dots$, получаем:

\begin{corollary}\label{l7:hensel-count}
Если $f(x_1) \equiv 0 \Mod p$ и $f'(x_1) \not\equiv 0 \Mod p$ ($x_1$ "--- \emph{простой}
корень), то для каждого $k$ существует единственный корень по модулю $p^k$, сравнимый с $x_1$ по
модулю $p$. Если все корни $f$ по модулю $p$ простые, то число корней по модулю $p^k$ равно
числу корней по модулю $p$.
\end{corollary}

\begin{algorithm}[подъём Гензеля]
\begin{algorithmic}[1]
\Require $f \in \ZZ[X]$, простой корень $x_1$ по модулю $p$, $k \ge 1$.
\Ensure корень $x$ по модулю $p^k$, $x \equiv x_1 \Mod p$.
\State $c \gets f'(x_1)^{-1} \bmod p$, $x \gets x_1$
\For{$j = 1, \dots, k - 1$}
  \State $s \gets \bigl(f(x) \bmod p^{j+1}\bigr) / p^j$ \Comment{деление нацело}
  \State $x \gets x - s c\, p^j \bmod p^{j+1}$
\EndFor
\State \Return $x$
\end{algorithmic}
\end{algorithm}

\begin{remark}[квадратичный подъём, метод Ньютона]\label{l7:newton}
Если $f(x) \equiv 0 \Mod{p^k}$ и $f'(x)$ обратим по модулю $p$, то
$x' = x - f(x)\, u \bmod p^{2k}$, где $u \equiv f'(x)^{-1} \Mod{p^{2k}}$, "--- корень по модулю
$p^{2k}$. Действительно, $h = -f(x)u$ делится на $p^k$, поэтому $h^2 \equiv 0 \Mod{p^{2k}}$ и по
лемме~\ref{l7:taylor} $f(x + h) \equiv f(x)(1 - u f'(x)) \equiv 0 \Mod{p^{2k}}$. Точность удваивается
на каждом шаге, и для модуля $p^k$ достаточно $\lceil \log_2 k \rceil$ шагов (вместо $k - 1$).
Так же устроены быстрое вычисление обратного по модулю $2^k$ для умножения Монтгомери
(лекция~5 первого семестра) и деление через умножение на «обратное» (лекция~3 первого
семестра).
\end{remark}

\subsection{\texorpdfstring{Квадратичные сравнения по модулю $p^k$}{Квадратичные сравнения по модулю p\textasciicircum k}}

Применим лемму Гензеля к $f(x) = x^2 - a$, $f'(x) = 2x$.

\begin{theorem}\label{l7:sqrt-pk}
Пусть $p$ "--- нечётное простое, $p \nmid a$, $k \ge 1$. Сравнение $x^2 \equiv a \Mod{p^k}$ имеет
$1 + \legendre{a}{p}$ решений, т.\,е. два решения $\pm x$, если $a$ "--- квадратичный вычет по
модулю $p$, и ни одного иначе. Корень получается из корня $x_1$ по модулю $p$ подъёмом
\[
  x_{j+1} = x_j + t_j p^j, \qquad
  t_j \equiv -\frac{x_j^2 - a}{p^j}\,(2x_1)^{-1} \Mod p .
\]
\end{theorem}

\begin{proof}
Всякое решение по модулю $p^k$ является решением и по модулю $p$; поэтому при
$\legendre{a}{p} = -1$ решений нет. Пусть $\legendre{a}{p} = 1$ и $\pm x_1$ "--- два корня по модулю
$p$. Так как $p$ нечётно и $p \nmid x_1$, $f'(\pm x_1) = \pm 2x_1 \not\equiv 0 \Mod p$: оба
корня простые. По следствию~\ref{l7:hensel-count} по модулю $p^k$ ровно два корня, и если $x$ "---
один из них, то $-x$ "--- другой.
\end{proof}

\begin{example}\label{l7:ex-hensel}
$x^2 \equiv 2 \Mod{7^3 = 343}$. По модулю $7$: $3^2 = 9 \equiv 2$, $x_1 = 3$,
$(2x_1)^{-1} = 6^{-1} \equiv 6 \Mod 7$.
\[
  \begin{array}{c|r|r|r|r|r}
    j & x_j & x_j^2 - 2 \bmod 7^{j+1} & s = (x_j^2 - 2)/7^j & t_j = -6s \bmod 7 & x_{j+1} \\ \hline
    1 & 3 & 7 & 1 & 1 & 10 \\
    2 & 10 & 98 & 2 & 2 & 108
  \end{array}
\]
Проверка: $108^2 = 11664 = 34 \cdot 343 + 2$. Ответ: $x \equiv \pm 108 \Mod{343}$, т.\,е.
$\{108, 235\}$.
\end{example}

\taskexample{sqrt-hensel}

\begin{remark}[случай $p \mid a$]
Если $a = p^{e} b$, $p \nmid b$, $e < k$, то решения есть только при чётном $e = 2c$ и
$\legendre{b}{p} = 1$, и тогда $x = p^{c} y$, где $y^2 \equiv b \Mod{p^{k - 2c}}$; при этом $y$
определён лишь по модулю $p^{k-2c}$, а $x$ "--- по модулю $p^{k - c}$, так что решений по модулю
$p^k$ становится больше двух ($2p^{c}$ штук). Например, $x^2 \equiv 9 \Mod{81}$ имеет $2 \cdot 3 = 6$ решений
$x = 3y$, $y \equiv \pm 1 \Mod 9$. Причина "--- кратный корень: по модулю $p$ сравнение
превращается в $x^2 \equiv 0$, $f'(0) = 0$, и лемма Гензеля неприменима.
\end{remark}

\begin{remark}[модуль $2^k$]\label{l7:two}
При $p = 2$ производная $2x$ всегда чётна, и лемма Гензеля в простой форме неприменима. Ответ
таков: для нечётного $a$ сравнение $x^2 \equiv a \Mod{2^k}$ имеет $1$ решение при $k = 1$;
при $k = 2$ "--- $2$ решения, если $a \equiv 1 \Mod 4$, и $0$ иначе; при $k \ge 3$ "--- $4$ решения
$\pm x, \pm x + 2^{k-1}$, если $a \equiv 1 \Mod 8$, и $0$ иначе. Например, $x^2 \equiv 1 \Mod 8$:
$x \in \{1, 3, 5, 7\}$; $x^2 \equiv 9 \Mod{16}$: $x \in \{3, 5, 11, 13\}$.

Набросок доказательства: квадрат нечётного числа $\equiv 1 \Mod 8$, откуда необходимость. Если
$x^2 \equiv a \Mod{2^k}$, $k \ge 3$, то $(x + 2^{k-1})^2 = x^2 + 2^k x + 2^{2k-2} \equiv x^2 + 2^k
\Mod{2^{k+1}}$, поэтому одно из чисел $x$, $x + 2^{k-1}$ "--- корень по модулю $2^{k+1}$
(модифицированный подъём). Если $x^2 \equiv y^2 \Mod{2^k}$, то $(x - y)(x + y) \equiv 0$, оба
множителя чётны, а их сумма $2x \equiv 2 \Mod 4$, так что один из них $\equiv 2 \Mod 4$, а другой
делится на $2^{k-1}$: $y \equiv \pm x \Mod{2^{k-1}}$, т.\,е. ровно четыре корня.
\end{remark}

\checkyourself{Почему при подъёме корня сравнения $x^2 \equiv a \Mod{p^k}$ достаточно один раз
найти $(2x_1)^{-1}$ по модулю $p$, а не обращать $2x_j$ по модулю $p^{j+1}$ на каждом шаге?}

\checkyourself{Сколько решений имеет сравнение $x^2 \equiv 1 \Mod{2^{10}}$? А
$x^2 \equiv 1 \Mod{3^{10}}$?}

\section{Составной модуль}\label{l7:sec-composite}

Пусть $n = m_1 m_2 \cdots m_r$, где $m_i = p_i^{e_i}$ "--- степени различных простых. По китайской
теореме об остатках (лекция~2) $\Zn{n} \cong \Zn{m_1} \times \cdots \times
\Zn{m_r}$, поэтому $x^2 \equiv a \Mod n$ равносильно системе $x^2 \equiv a \Mod{m_i}$,
$i = 1, \dots, r$, и каждый набор корней $(x \bmod m_1, \dots, x \bmod m_r)$ даёт ровно один корень
по модулю $n$.

\begin{theorem}\label{l7:crt-count}
Пусть $n$ нечётно, $n = p_1^{e_1} \cdots p_r^{e_r}$, $\gcd(a, n) = 1$. Число решений сравнения
$x^2 \equiv a \Mod n$ равно
\[
  \prod_{i=1}^{r} \Bigl(1 + \legendre{a}{p_i}\Bigr) \in \{0, 2^r\}.
\]
Решения есть тогда и только тогда, когда $a$ "--- квадратичный вычет по модулю каждого $p_i$.
\end{theorem}

\begin{proof}
По КТО число решений равно произведению чисел решений по модулям $p_i^{e_i}$, а они даются
теоремой~\ref{l7:sqrt-pk}.
\end{proof}

Если $n$ чётно, к произведению добавляется множитель $1$, $2$, $4$ или $0$ для степени двойки
(замечание~\ref{l7:two}). Если $p_i \mid a$ для некоторых $i$ (а $n$ свободно от квадратов), то по
модулю $p_i$ корень один ($x \equiv 0$), и множитель $1 + \legendre{a}{p_i}$ равен $1$.

\emph{Как вычислять.} Удобно один раз построить «базис» КТО: $e_i = M_i y_i \bmod n$, где
$M_i = n/m_i$, $y_i = M_i^{-1} \bmod m_i$, так что $e_i \equiv 1 \Mod{m_i}$ и $e_i \equiv 0$ по
остальным модулям. Тогда корень с остатками $(r_1, \dots, r_r)$ равен $\sum r_i e_i \bmod n$, и
все $2^r$ корней $\sum (\pm r_i) e_i$ получаются без повторного решения систем.

\begin{example}\label{l7:ex-77}
$x^2 \equiv 23 \Mod{77}$, $77 = 7 \cdot 11$. По модулю $7$: $23 \equiv 2$, $7 \equiv 3 \Mod 4$,
$2^{(7+1)/4} = 4$, $4^2 = 16 \equiv 2$: корни $\pm 4$. По модулю $11$: $23 \equiv 1$, корни $\pm 1$.
Базис: $e_1 = 11 \cdot (11^{-1} \bmod 7) = 11 \cdot 2 = 22$, $e_2 = 7 \cdot (7^{-1} \bmod 11) =
7 \cdot 8 = 56$. Корни $\pm 4 \cdot 22 \pm 1 \cdot 56$:
\[
  88 + 56 = 144 \equiv 67, \quad 88 - 56 = 32, \quad -32 \equiv 45, \quad -67 \equiv 10.
\]
Ответ: $x \in \{10, 32, 45, 67\}$.
\end{example}

\taskexample{sqrt-composite}

\checkyourself{Сколько решений имеет сравнение $x^2 \equiv 1 \Mod{105}$? Перечислите их, не решая
восемь систем.}

\section{\texorpdfstring{Корни по модулю $N = pq$ и факторизация}{Корни по модулю N = pq и факторизация}}\label{l7:sec-rabin}

Для простого $p$ корни извлекаются за полиномиальное время, и по теореме~\ref{l7:crt-count} то же
верно для модуля $N = pq$, \emph{если известно разложение} $N$. Обратно, умение извлекать корни
по модулю $N$ позволяет разложить $N$.

\begin{theorem}[Рабин]\label{l7:rabin}
Пусть $N = pq$, $p \ne q$ "--- нечётные простые.
\begin{enumerate}
  \item Если известны два корня $x, y$ сравнения $x^2 \equiv y^2 \equiv a \Mod N$, $\gcd(a, N) = 1$,
    с $y \not\equiv \pm x \Mod N$, то $\gcd(x - y, N) \in \{p, q\}$.
  \item Если есть алгоритм $\mathcal{A}$, который по любому квадратичному вычету
    $a \in \Znx{N}$ возвращает какой-нибудь его квадратный корень, то $N$ раскладывается
    вероятностным алгоритмом, делающим в среднем не более двух обращений к $\mathcal{A}$ и
    $\bigO(1)$ вычислений НОД.
\end{enumerate}
\end{theorem}

\begin{proof}
1. Из $x^2 \equiv y^2$ следует $N \mid (x - y)(x + y)$. Если бы $N \mid x - y$ или $N \mid x + y$,
было бы $y \equiv \pm x$. Значит, $p$ и $q$ делят разные множители (каждый делит хотя бы один, так как они простые), и
$\gcd(x - y, N)$ "--- это $p$ или $q$.

2. Выберем случайное $x \in \{1, \dots, N - 1\}$. Если $\gcd(x, N) > 1$, разложение найдено.
Иначе вычислим $a = x^2 \bmod N$ и $y = \mathcal{A}(a)$. По теореме~\ref{l7:crt-count} у $a$ ровно
четыре корня: $\pm x$ и $\pm x'$, где $x' \equiv x \Mod p$, $x' \equiv -x \Mod q$. Все четыре
корня дают одно и то же $a$, поэтому при фиксированном $a$ случайное $x$ равновероятно
является любым из них, а ответ $y$ от $x$ не зависит (алгоритм видит только $a$). Значит, с
вероятностью $\frac12$ выполнено $y \not\equiv \pm x$, и по п.~1 $\gcd(x - y, N)$ "--- нетривиальный
делитель. Повторяя, получаем в среднем не более двух попыток.
\end{proof}

Например, для $N = 77$ из примера~\ref{l7:ex-77}: корни $10$ и $32$ числа $23$ не отличаются
знаком, и $\gcd(32 - 10, 77) = \gcd(22, 77) = 11$, $\gcd(32 + 10, 77) = 7$.

\begin{remark}[криптосистема Рабина]
Открытый ключ "--- $N = pq$ (обычно $p \equiv q \equiv 3 \Mod 4$), секретный "--- $p, q$.
Шифрование: $c = m^2 \bmod N$. Расшифрование: $m_p = c^{(p+1)/4} \bmod p$,
$m_q = c^{(q+1)/4} \bmod q$ (теорема~\ref{l7:p34}), затем по КТО четыре кандидата
$(\pm m_p, \pm m_q)$; верный выбирают по заранее оговорённой избыточности открытого текста
(например, повторённым последним битам). По теореме~\ref{l7:rabin} расшифрование без знания $p, q$
\emph{равносильно} факторизации $N$ "--- в отличие от RSA, для которого такая эквивалентность не
доказана~\cite{rabin-signatures}. Обратная сторона: противник, получивший доступ к
расшифрованию (атака с выбором шифртекста), отправляет $c = x^2$ и с вероятностью $\frac12$
разлагает $N$ по п.~1 теоремы. Для $N = 77$ и $c = 23$ кандидаты "--- $10, 32, 45, 67$.
\end{remark}

\checkyourself{Почему в теореме Рабина важно, что алгоритм $\mathcal{A}$ получает только $a$, а не $x$?
Что можно сказать о корнях $x$ и $x'$ из доказательства, если вычислить $\gcd(x - x', N)$?}

\checkyourselfsection

\begin{furtherreading}
  \readitem[разд.~2.3--2.4]{glukhov}{квадратичные вычеты, критерий Эйлера, метод Чиполлы;
    сравнения по модулю степени простого}
  \readitem[гл.~IV--V]{vinogradov}{сравнения по составному модулю и модулю $p^\alpha$;
    символы Лежандра и Якоби, закон взаимности}
  \readitem[разд.~2.3]{crandall-pomerance}{квадратичные вычеты и алгоритмы извлечения квадратных
    корней (Тоннелли, Аткин, Чиполла), корни по модулю степени простого}
  \readitem[разд.~1.5]{cohen}{символ Лежандра"--~Якоби, алгоритм Тоннелли"--~Шенкса}
  \readitem[гл.~12]{shoup}{закон взаимности и вычисление квадратных корней по модулю}
  \readitem[разд.~3.9]{hps}{квадратичные вычеты и закон взаимности в криптографическом контексте}
  \readitem[гл.~5]{ireland-rosen}{доказательства свойств символа Лежандра и закона взаимности}
  \readitem{rabin-signatures}{криптосистема и подпись Рабина, эквивалентность факторизации}
  \readitem{msu-ent}{курс элементарной теории чисел: квадратичные вычеты, лемма Гензеля}
  \readitem{wiki-cipolla}{алгоритм Чиполлы с примерами и обобщением на степени простого}
\end{furtherreading}
\end{lecture}
